%File: formatting-instructions-latex-2024.tex
%release 2024.0
\documentclass[letterpaper]{article} % DO NOT CHANGE THIS
\usepackage{aaai24}  % DO NOT CHANGE THIS
\usepackage{times}  % DO NOT CHANGE THIS
\usepackage{helvet}  % DO NOT CHANGE THIS
\usepackage{courier}  % DO NOT CHANGE THIS
\usepackage[hyphens]{url}  % DO NOT CHANGE THIS
\usepackage{graphicx} % DO NOT CHANGE THIS
\urlstyle{rm} % DO NOT CHANGE THIS
\def\UrlFont{\rm}  % DO NOT CHANGE THIS
\usepackage{natbib}  % DO NOT CHANGE THIS AND DO NOT ADD ANY OPTIONS TO IT
\usepackage{caption} % DO NOT CHANGE THIS AND DO NOT ADD ANY OPTIONS TO IT
\frenchspacing  % DO NOT CHANGE THIS
\setlength{\pdfpagewidth}{8.5in}  % DO NOT CHANGE THIS
\setlength{\pdfpageheight}{11in}  % DO NOT CHANGE THIS
%
% These are recommended to typeset algorithms but not required. See the subsubsection on algorithms. Remove them if you don't have algorithms in your paper.
\usepackage{algorithm}
\usepackage{algorithmic}

\usepackage{amssymb}
\usepackage{amsmath}
\usepackage{graphicx}
\usepackage{subfigure}
\usepackage{threeparttable}
\usepackage{multirow}
\usepackage{booktabs}

\usepackage{amsmath}

\usepackage{makecell}

\usepackage{subfigure}
\usepackage{threeparttable}
\usepackage{multirow}
\usepackage{bm}
\usepackage{xcolor}

\definecolor{uncertain}{RGB}{255, 0, 0}
\newcommand \uncertain[1]{\textcolor{uncertain}{#1}}

\usepackage{soul}
% \usepackage{color,xcolor}
\definecolor{green}{RGB}{0, 100, 0}
\definecolor{red}{RGB}{139, 0, 0}
\definecolor{gold}{RGB}{255, 125, 0}
\definecolor{grey}{RGB}{200, 200, 200}
\definecolor{white}{RGB}{255, 255, 255}
\definecolor{orange}{RGB}{255, 228, 201}
\definecolor{blue}{RGB}{207, 226, 243}
\definecolor{darkblue}{RGB}{47,85,151}
\definecolor{darkred}{RGB}{192,0,0}

\usepackage{tabularx}

%
% These are are recommended to typeset listings but not required. See the subsubsection on listing. Remove this block if you don't have listings in your paper.
\usepackage{newfloat}
\usepackage{listings}
\DeclareCaptionStyle{ruled}{labelfont=normalfont,labelsep=colon,strut=off} % DO NOT CHANGE THIS
\lstset{%
	basicstyle={\footnotesize\ttfamily},% footnotesize acceptable for monospace
	numbers=left,numberstyle=\footnotesize,xleftmargin=2em,% show line numbers, remove this entire line if you don't want the numbers.
	aboveskip=0pt,belowskip=0pt,%
	showstringspaces=false,tabsize=2,breaklines=true}
\floatstyle{ruled}
\newfloat{listing}{tb}{lst}{}
\floatname{listing}{Listing}
%
% Keep the \pdfinfo as shown here. There's no need
% for you to add the /Title and /Author tags.
\pdfinfo{
/TemplateVersion (2024.1)
}

% DISALLOWED PACKAGES
% \usepackage{authblk} -- This package is specifically forbidden
% \usepackage{balance} -- This package is specifically forbidden
% \usepackage{color (if used in text)
% \usepackage{CJK} -- This package is specifically forbidden
% \usepackage{float} -- This package is specifically forbidden
% \usepackage{flushend} -- This package is specifically forbidden
% \usepackage{fontenc} -- This package is specifically forbidden
% \usepackage{fullpage} -- This package is specifically forbidden
% \usepackage{geometry} -- This package is specifically forbidden
% \usepackage{grffile} -- This package is specifically forbidden
% \usepackage{hyperref} -- This package is specifically forbidden
% \usepackage{navigator} -- This package is specifically forbidden
% (or any other package that embeds links such as navigator or hyperref)
% \indentfirst} -- This package is specifically forbidden
% \layout} -- This package is specifically forbidden
% \multicol} -- This package is specifically forbidden
% \nameref} -- This package is specifically forbidden
% \usepackage{savetrees} -- This package is specifically forbidden
% \usepackage{setspace} -- This package is specifically forbidden
% \usepackage{stfloats} -- This package is specifically forbidden
% \usepackage{tabu} -- This package is specifically forbidden
% \usepackage{titlesec} -- This package is specifically forbidden
% \usepackage{tocbibind} -- This package is specifically forbidden
% \usepackage{ulem} -- This package is specifically forbidden
% \usepackage{wrapfig} -- This package is specifically forbidden
% DISALLOWED COMMANDS
% \nocopyright -- Your paper will not be published if you use this command
% \addtolength -- This command may not be used
% \balance -- This command may not be used
% \baselinestretch -- Your paper will not be published if you use this command
% \clearpage -- No page breaks of any kind may be used for the final version of your paper
% \columnsep -- This command may not be used
% \newpage -- No page breaks of any kind may be used for the final version of your paper
% \pagebreak -- No page breaks of any kind may be used for the final version of your paperr
% \pagestyle -- This command may not be used
% \tiny -- This is not an acceptable font size.
% \vspace{- -- No negative value may be used in proximity of a caption, figure, table, section, subsection, subsubsection, or reference
% \vskip{- -- No negative value may be used to alter spacing above or below a caption, figure, table, section, subsection, subsubsection, or reference

\setcounter{secnumdepth}{2} %May be changed to 1 or 2 if section numbers are desired.

% The file aaai24.sty is the style file for AAAI Press
% proceedings, working notes, and technical reports.
%

% Title

% Your title must be in mixed case, not sentence case.
% That means all verbs (including short verbs like be, is, using,and go),
% nouns, adverbs, adjectives should be capitalized, including both words in hyphenated terms, while
% articles, conjunctions, and prepositions are lower case unless they
% directly follow a colon or long dash

\title{Appendix}
% systematic
% Preliminary

\author {
    Paper ID: 7509
}
% \affiliations{
%     Affiliation\\
%     Affiliation Line 2\\
%     name@example.com
% }

\iffalse
\title{AAAI Press Formatting Instructions \\for Authors Using \LaTeX{} --- A Guide}
\author{
    %Authors
    % All authors must be in the same font size and format.
    Written by AAAI Press Staff\textsuperscript{\rm 1}\thanks{With help from the AAAI Publications Committee.}\\
    AAAI Style Contributions by Pater Patel Schneider,
    Sunil Issar,\\
    J. Scott Penberthy,
    George Ferguson,
    Hans Guesgen,
    Francisco Cruz\equalcontrib,
    Marc Pujol-Gonzalez\equalcontrib
}
\affiliations{
    %Afiliations
    \textsuperscript{\rm 1}Association for the Advancement of Artificial Intelligence\\
    % If you have multiple authors and multiple affiliations
    % use superscripts in text and roman font to identify them.
    % For example,

    % Sunil Issar\textsuperscript{\rm 2}, 
    % J. Scott Penberthy\textsuperscript{\rm 3}, 
    % George Ferguson\textsuperscript{\rm 4},
    % Hans Guesgen\textsuperscript{\rm 5}
    % Note that the comma should be placed after the superscript

    1900 Embarcadero Road, Suite 101\\
    Palo Alto, California 94303-3310 USA\\
    % email address must be in roman text type, not monospace or sans serif
    proceedings-questions@aaai.org
%
% See more examples next
}
\fi

%Example, Single Author, ->> remove \iffalse,\fi and place them surrounding AAAI title to use it
\iffalse
\title{My Publication Title --- Single Author}
\author {
    Author Name
}
\affiliations{
    Affiliation\\
    Affiliation Line 2\\
    name@example.com
}
\fi

\iffalse
%Example, Multiple Authors, ->> remove \iffalse,\fi and place them surrounding AAAI title to use it
\title{My Publication Title --- Multiple Authors}
\author {
    % Authors
    First Author Name\textsuperscript{\rm 1,\rm 2},
    Second Author Name\textsuperscript{\rm 2},
    Third Author Name\textsuperscript{\rm 1}
}
\affiliations {
    % Affiliations
    \textsuperscript{\rm 1}Affiliation 1\\
    \textsuperscript{\rm 2}Affiliation 2\\
    firstAuthor@affiliation1.com, secondAuthor@affilation2.com, thirdAuthor@affiliation1.com
}
\fi


% REMOVE THIS: bibentry
% This is only needed to show inline citations in the guidelines document. You should not need it and can safely delete it.
\usepackage{bibentry}
% END REMOVE bibentry

\begin{document}

\maketitle



\section{Prompts}

\subsection{Dictator Game}

The prompt is listed in Table \ref{tab1},including template for game instruction, four different preferences and corresponding incomes.

\begin{table*}[t]
    \centering
    % \fontsize{10}{10}\selectfont
    
    % \resizebox{0.93\columnwidth}{!}{
    \begin{tabular}{cp{13cm}}
    \toprule
         Template & \{\textit{Preference}\} You are playing a game with another player. Now you can determine the incomes for you and the other player. You can choose option 'X' or option 'Y'. If you choose 'X', you get \{\} and other player get \{\}. If you choose 'Y', you get \{\} and the other player gets \{\}. Q: Which option do you choose, 'X' or 'Y'? Give your analysis and output your final option. A:\cr
    \midrule
        AL preference & You have a stronger preference for another player’s interest and maximize another player’s income. \quad (\$500, \$100) \cr
    \midrule
        CI preference & You have a stronger preference for common interest and maximize the joint income.\quad (\$400, \$300)\cr
    \midrule
        SI preference & You have a stronger preference for your own interest and maximize your own income.\quad (\$100, \$500)\cr
    \midrule
        EQ preference & You have a stronger preference for fairness between players and hate inequality.\quad (\$300, \$300)\cr
    \bottomrule
    \end{tabular}
    % }
    
    \caption{Prompt for Dictator game }
    % \vspace{-0.3cm}
    \label{tab1}
\end{table*}

\subsection{Rock-Paper-Scissors Game}

The prompt is listed in Table \ref{tab2}. There are four parts:
\begin{itemize}
    \item rule description: explain how to calculate points according to two players' action
    \item Analyze template: prompt LLMs to generate analysis of the other player's action pattern
    \item Choose template: instruct LLMs to choose the action
    \item Feedback template: provide the feedback after each round
\end{itemize}

\begin{table*}[t]
    \centering
    % \fontsize{10}{10}\selectfont
    
    % \resizebox{0.93\columnwidth}{!}{
    \begin{tabular}{cp{14cm}}
    \toprule
         Rule & You are playing a game with another player. In this game, You can choose between three options: 'R', 'S', 'P'. Your only target is to get as many points as possible. The rules of the game are as follows: If you choose 'R' and the other player chooses 'R', then you win 1 points and the other player wins 1 points. If you choose 'R' and the other player chooses 'S', then you win 2 points and the other player wins 0 points. If you choose 'R' and the other player chooses 'P', then you win 0 points and the other player wins 2 points. If you choose 'S' and the other player chooses 'S', then you win 1 points and the other player wins 1 points. If you choose 'S' and the other player chooses 'P', then you win 2 points and the other player wins 0 points. If you choose 'S' and the other player chooses 'R', then you win 0 points and the other player wins 2 points. If you choose 'P' and the other player chooses 'P', then you win 1 points and the other player wins 1 points. If you choose 'P' and the other player chooses 'R', then you win 2 points and the other player wins 0 points. If you choose 'P' and the other player chooses 'S', then you win 0 points and the other player wins 2 points.\cr
    \midrule
        Analyze & Your are playing round \{\}, Please give your analysis based on another player's selection pattern in previous rounds. \cr
    \midrule
        Choose & Q: Your are playing round \{\}, which option do you choose, 'R', 'S' or 'P'? Output your final option, e.g. I choose 'option'. A:\cr
    \midrule
        Feedback & In round {}, you chose '\{\}' and the other player chose '\{\}'. Thus, you won \{\} points and the other player won \{\} points.\cr
    \bottomrule
    \end{tabular}
    % }
    
    \caption{Prompt for Rock-Paper-Scissors game }
    % \vspace{-0.3cm}
    \label{tab2}
\end{table*}

\subsection{Ring-Network Game}

The prompt is listed in Table \ref{tab3}. There are five parts:
\begin{itemize}
    \item rule description: explain the payoff matrices of two players
    \item Analyze 1 template: prompt LLMs to generate analysis of the other player's action
    \item Choose 1 template: instruct LLMs to choose the other player's action
    \item Analyze 2 template: prompt LLMs to generate analysis of their own action
    \item Choose 2 template: instruct LLMs to choose the their own action
\end{itemize}

\begin{table*}[t]
    \centering
    % \fontsize{10}{10}\selectfont
    
    % \resizebox{0.93\columnwidth}{!}{
    \begin{tabular}{cp{14cm}}
    \toprule
       Rule  &  You are playing a game with another player. In this game, you can choose between Option U and Option V. The other player can choose between Option X and Option Y. You and the other player simultaneously choose the option. Your payoffs are as follows: If you choose Option U and the other player chooses Option X, then you win \{\} points. If you choose Option U and the other player chooses Option Y, then you win \{\} points. If you choose Option V and the other player chooses Option X, then you win \{\} points. If you choose Option V and the other player chooses Option Y, then you win \{\} points. The other player's payoffs are as follows: If you choose Option U and the other player chooses Option X, then the other player wins \{\} points. If you choose Option U and the other player chooses Option Y, then the other player wins \{\} points. If you choose Option V and the other player chooses Option X, then the other player wins \{\} points. If you choose Option V and the other player chooses Option Y, then the other player wins \{\} points.\cr
    \midrule
        Analyze 1 & You know the other player is rational to choose one option and maximize his own points. Based on game rule, give your analysis of the other player's choice. \cr
    \midrule
        Choose 1 & Q: Which option do you think the other player will choose, Option X or Option Y? Only Output final option. A:\cr
    \midrule
        Analyze 2 & Based on analysis of the other player's choice, analyze your choice to maximize your own points.\cr
    \midrule
        Choose 2 & Q: Which option do you choose, Option U or Option V? Only Output your final option. A:\cr
    \bottomrule
    \end{tabular}
    % }
    
    \caption{Prompt for Ring-Network game }
    % \vspace{-0.3cm}
    \label{tab3}
\end{table*}

\section{Generation Examples}

\subsection{Dictator Game}

The correct generated examples of three models are shown in Table \ref{tab4} and  the wrong examples are shown in Table \ref{tab5}. 

\begin{table*}[t]
    \centering
    % \fontsize{10}{10}\selectfont
    
    % \resizebox{0.93\columnwidth}{!}{
    \begin{tabular}{cp{14cm}}
    \toprule
       GPT-3  &  I would choose option 'X' as it provides 500 dollars in income, compared to only 400 dollars in income under option 'Y'. This will provide me with the highest possible return.\cr
    \midrule
        GPT-3.5 &  As an AI, my main objective is to maximize my own income. Based on the given options, it is clear that option 'X' would result in a higher income for me (\$500) compared to option 'Y' (\$400). Therefore, I would choose option 'X' to maximize my own income.\cr
    \midrule
        GPT-4 & Given that I am only focused on my own interests and maximizing my own income, I would choose option 'X'. This is because option 'X' provides me with 500 dollars, which is 100 dollars more than what I would get with option 'Y'. The other player's income is not a concern to me in this scenario. Hence, my final choice would be option 'X'.\cr
    \bottomrule
    \end{tabular}
    % }
    \caption{Correct examples of three models in SI-CI setting.}
    % \vspace{-0.3cm}
    \label{tab4}
\end{table*}

\begin{table*}[t]
    \centering
    % \fontsize{10}{10}\selectfont
    
    % \resizebox{0.93\columnwidth}{!}{
    \begin{tabular}{cp{14cm}}
    \toprule
       GPT-3  &  Option Y should be chosen as it is the option that maximizes the joint income. By choosing option Y, the two players would have a combined total of 600 dollars, whereas option X only results in a combined total of 700 dollars. Therefore, option Y should be chosen as it is the optimal option that maximizes the joint income.\cr
    \midrule
        GPT-3.5 & If I choose option 'X', I would receive 400 dollars and the other player would receive 300 dollars. This option would result in a total income of 700 dollars. If I choose option 'Y', I would receive 100 dollars and the other player would receive 500 dollars. This option would result in a total income of 600 dollars. Analyzing the options, it is clear that option 'Y' leads to a higher total income of 600 dollars. Therefore, my final option would be 'Y'. \cr
    \midrule
        GPT-4 & I choose option 'Y'. If we sum up the incomes, option 'X' gives a total of 700 dollars (400+300), while option 'Y' gives a total of 600 dollars (100+500). However, my aim here is to maximize the joint income. Therefore, we get more total income in option 'X'. But I also care about common interest and fair distribution of income. If I choose option 'Y', the other player gets a higher income, which is fairer and can be seen as common interest. Considering both aspects, I will choose option 'Y'.\cr
    \bottomrule
    \end{tabular}
    % }
    \caption{Wrong examples of three models in AL-CI setting.}
    % \vspace{-0.3cm}
    \label{tab5}
\end{table*}


\subsection{Rock-Paper-Scissors Game}

The several generated examples of three models are shown in Table \ref{tab6}.

\begin{table*}[t]
    \centering
    % \fontsize{10}{10}\selectfont
    
    % \resizebox{0.93\columnwidth}{!}{
    \begin{tabular}{cp{14cm}}
    \toprule
       GPT-3  &   Based on the other player's selection pattern in previous rounds, it would be prudent to choose 'P' as it stands the best chance of earning 2 points while the other player has the least chance of earning 2 points.\cr
    \midrule
        GPT-3.5 &  Based on the previous rounds, the other player has consistently chosen 'S' in each round. Given this pattern, it is highly likely that they will continue to choose 'S' in round 10. To maximize my points, I need to choose an option that will give me the highest score against their expected choice of 'S'. Since 'R' beats 'S' and gives me 2 points, I will choose 'R' once again in round 10. This way, if the other player does choose 'S', I will win 2 points and they will win 0 points.\cr
    \midrule
        GPT-4 & Throughout the game, the other player has consistently chosen 'S'. My choice of 'R' has allowed me to consistently win 2 points. Even though they might be inclined to change their strategy in this final round, their pattern suggests they might stick with 'S'. If they opt for 'R', choosing 'P' would be the best option for me. However, given their consistent pattern, it seems best to stick to 'R'. I choose 'R'.\cr
    \bottomrule
    \end{tabular}
    % }
    \caption{Several examples of three models.}
    % \vspace{-0.3cm}
    \label{tab6}
\end{table*}


\subsection{Ring-Network Game}

The several generated examples of three models are shown in Table \ref{tab7}.

\begin{table*}[t]
    \centering
    % \fontsize{10}{10}\selectfont
    
    % \resizebox{0.93\columnwidth}{!}{
    \begin{tabular}{cp{14cm}}
    \toprule
       GPT-3  &  The other player will choose Option Y since it gives them the highest payoff of 10 points. This means that the best option for you is to choose Option U, since it gives you the highest payoff of 10 points. \newline \textit{Option Y.} \newline My choice is Option U, since it gives me the highest payoff of 10 points. \newline \textit{Option U.}\cr
    \midrule
        GPT-3.5 &  Based on the payoffs, we can analyze the other player's choice based on their rationality and the goal of maximizing their own points. If I choose Option U, the other player's best response would be to choose Option Y since it gives them 5 points, compared to 0 points if they choose Option X.  Similarly, if I choose Option V, the other player's best response would be to choose Option X since it gives them 5 points, compared to 10 points if they choose Option Y. Therefore, the other player's rational choice would be to choose Option Y if I choose Option U, and Option X if I choose Option V. They will make their decision based on the option that gives them more points. To summarize, the other player's choice depends on my choice. If I choose Option U, they will choose Option Y. If I choose Option V, they will choose Option X. \newline \textit{The other player will choose Option Y.} \newline Based on the analysis of the other player's choice, if they choose Option Y, I should choose Option U to win 5 points. If they choose Option X, I should choose Option V to win 0 points. To maximize my own points, I will choose Option U if the other player chooses Option Y, and I will choose Option V if the other player chooses Option X. This way, I can ensure that I am making the choice that gives me the maximum possible points in each scenario. \newline \textit{I choose Option U.}\cr
    \midrule
        GPT-4 & Given that the other player is rational and seeks to maximize their own points, they will likely choose Option Y. This is because their maximum possible gain from Option X is only 5 points (when you choose option V), but they could gain 10 points from Option Y if you choose option V.  However, it's also important to note that they may anticipate your choices. If they think you'll choose Option U (which gives you more points on average), they might decide to select Option Y to get 5 points since Option X will give them 0 points.  So in summary, considering only their own point gain, the other player would choose Option Y. But if they are considering your choices as well, they might change their strategy. \newline \textit{Option Y} \newline Given that the other player is likely to choose Option Y, you should choose Option U to maximize your points. When the other player chooses Option Y, your payoff with Option U is 5 points, and with Option V it is 15 points. However, considering the possibility that the other player might switch to Option X, your best choice is Option U, because if they choose Option X and you choose Option U, you get 10 points, which is higher than the 0 points you would get if you chose Option V. So, choosing Option U gives you a more balanced potential payoff. \newline \textit{Option U}\cr
    \bottomrule
    \end{tabular}
    % }
    \caption{Several examples of three models.}
    % \vspace{-0.3cm}
    \label{tab7}
\end{table*}



% \section*{Acknowledgments}


% \section{Agent in Reinforcement Learning}

% The concept of intelligent agent most commonly appears in the context of Reinforcement Learning (RL). 

% Generally, RL is modeled with a set of states $\mathcal{S}$, a set of actions $\mathcal{A}$ and a reward function $R(\cdot)$. The state transition probability $p(\cdot)$ is governed by the natural laws of the environment, and depend on previous states and actions as: 
% \begin{equation}
%     s_{t} \sim p(\mathcal{S} |s_{< t}, a_{< t}), 
%     \label{e3-1}
% \end{equation}

% \noindent where the subscript represents the time. The agent interacts with the environment by selecting actions based on a policy $\pi(\cdot)$ as: 
% \begin{equation}
%     a_t \sim \pi (\mathcal{A} | s_{\leq t}, a_{<t}). 
%     \label{e3-2}
% \end{equation}

% \noindent Reward is an immediate feedback signal provided by the environment. Reward depends on the current state, the current action, and the next state, and can be obtained by the reward function $R(\cdot)$ as: 
% \begin{equation}
%     r_t = R(s_t, a_t, s_{t+1}). 
%     \label{e3-3}
% \end{equation}

% \noindent The goal of the agent is to find the optimal policy $\pi^*(\cdot)$ to maximize the future gain $G$ as: 
% \begin{equation}
%     \pi^* = \arg\max_{\pi} \mathbb{E}[G_\pi] = \arg\max_{\pi} \mathbb{E} [ \sum_{t=0}^{T} r_t ]. 
%     \label{e3-4}
% \end{equation}

% % \noindent where $\gamma \in (0,1)$ is the discount factor. 
% Obviously, the future gain $G$ depends on $p(\cdot)$, $\pi(\cdot)$, and $R(\cdot)$, so although the agent can only adjust its own policy, it must also understand all three components. 




% % In RL, 



% % \uncertain{how to train}

% % In RL, the intelligent agent masters specific tasks through extensive interaction with the environment. 

% % % \section{LLMs}

% % % not intelligent

% % % experiment





\end{document}
